\chapter{Firewall Access Controls}
\label{ch:3.4}
\chaptertopic{Topic 3.4}

\chapterquote{Security is a process, not a product.}{Bruce Schneier, \emph{Crypto-Gram} (2000)}

\section{Motivation}

Chapter~\ref{ch:3.3} divided the network into zones. This chapter is about the gatekeeper that
decides what may cross between them, or between the network and the outside world at all: the
firewall. Two questions organize everything here. First, how much does the firewall actually
remember about a conversation as it decides -- nothing at all, or a full sense of which connections
are legitimately open? Second, given a list of rules, in what order does the firewall apply them? And what happens when
that order is wrong, even though every individual rule looks correct on its own?

\begin{apalignment}{Unit 3: Securing Networks}{Topic 3.4 -- Protecting Networks: Firewalls}
\begin{apobjectives}
\apobjective{3.4.A}{Mitigate Risk}{Identify the kinds of network-based firewall -- stateless packet
filters, stateful inspection, application proxies, and next-generation platforms -- and know which
layers each inspects and which decisions each can and cannot make.}
\apobjective{3.4.B}{Mitigate Risk}{Explain how a firewall uses an access control list to permit or
deny traffic in either direction, including first-match evaluation, the load-bearing role of rule
order, and the default-deny rule at the end of the list.}
\apobjective{3.4.C}{Mitigate Risk}{Determine where firewalls belong in a topology -- the perimeter,
the screened subnet around public services, and internal boundaries between zones of differing
sensitivity -- so that a single compromise is contained.}
\apobjective{3.4.D}{Mitigate Risk}{Configure a firewall so the flow of traffic matches a stated
policy, and verify by tracing which rule actually applies to a given packet.}
\end{apobjectives}
\end{apalignment}

\section{Stateless and Stateful Firewalls}

Every packet carries a header -- source and destination IP, protocol, source and destination port
-- and a payload. The simplest possible firewall looks only at header fields, decides for each
packet independently, and moves on.

\begin{definition}{Stateless (packet-filter) firewall}
A firewall that applies a fixed rule set to each packet in isolation, with no memory of any
previous packet or of whether a connection was ever legitimately established.
\end{definition}

That blindness is both a stateless firewall's strength and its weakness. It is fast and cheap --
routers, switches, and early border firewalls all worked this way. But it cannot tell an answer to
a legitimate request apart from an uninvited packet arriving out of nowhere. A rule reading ``allow
TCP to port 80'' will forward \emph{any} packet matching that header, including one an attacker
crafted purely to look like part of an existing conversation.

\begin{example}[Why statelessness fails against a spoofed reset]
An internal network's stateless firewall allows inbound TCP on port 80. An attacker outside spoofs
the source address of an internal web server and sends a single TCP packet with the RST flag set,
destination port 80. The stateless firewall checks the header, sees ``TCP, port 80, allow,'' and
forwards it -- disrupting an active connection inside the network, even though no real conversation
with that attacker ever existed.
\end{example}

\begin{definition}{Stateful (stateful-inspection) firewall}
A firewall that maintains a connection-tracking table recording every active session's source and
destination IP, ports, protocol, and connection state, permitting return traffic automatically for
tracked sessions and dropping any packet claiming membership in a session that was never
established.
\end{definition}

A stateful firewall builds this table directly from the TCP three-way handshake: the first SYN
packet, once it passes the rule set, is logged as \texttt{SYN\_SENT}; the handshake's completion
advances the entry to \texttt{ESTABLISHED}. Return traffic -- the SYN-ACK, subsequent ACKs, the
data itself -- is recognized as belonging to a tracked session and permitted. The same spoofed RST
packet from the earlier example is checked against the connection table instead, finds no
established session with those endpoints, and is dropped -- regardless of how valid its header
would look to a stateless rule.

\begin{lstlisting}[language=Python, caption={Stateless versus stateful packet classification}]
def stateless_classify(packet, rules):
    """rules: ordered list of dicts like {"proto": str, "dst_port": int,
       "action": "allow"|"deny"}. Returns "allow" or "deny" using
       first-match, defaulting to deny."""
    for rule in rules:
        if rule["proto"] == packet["proto"] and rule["dst_port"] == packet["dst_port"]:
            return rule["action"]
    return "deny"

def stateful_classify(packet, connections, rules):
    """connections: dict mapping (src_ip, src_port, dst_ip, dst_port) ->
       state string. Mutates connections as the handshake progresses.
       Returns "allow" or "deny"."""
    key = (packet["src_ip"], packet["src_port"], packet["dst_ip"], packet["dst_port"])
    reverse_key = (packet["dst_ip"], packet["dst_port"], packet["src_ip"], packet["src_port"])

    if packet["flag"] == "SYN":
        if stateless_classify(packet, rules) == "deny":
            return "deny"
        connections[key] = "SYN_SENT"
        return "allow"
    if packet["flag"] == "ACK" and reverse_key in connections:
        connections[reverse_key] = "ESTABLISHED"
        return "allow"
    if key in connections or reverse_key in connections:
        return "allow"
    return "deny"
\end{lstlisting}

Stateful inspection is not the same as content inspection: it tracks \emph{connection state}, not
application-layer payload. Malware delivered inside an HTTPS stream still passes through a stateful
firewall on port 443 exactly as any other traffic on that port would. The firewall only confirms
the connection is legitimately established -- it does not look at what is inside it. Content
inspection is a separate, later capability, layered on top by application proxies and
next-generation firewalls. Stateless firewalls remain useful at high-throughput edges where raw
speed outweighs fine-grained awareness; stateful inspection is the standard everywhere connection
context actually matters, which is most internal and server-facing deployments.

\subsection{Application Proxies and Next-Generation Firewalls}

Stateless and stateful firewalls both make their decision from the same limited evidence: headers
and, for stateful inspection, connection state. Two further firewall categories reach higher up the
protocol stack, into the payload itself.

\begin{definition}{Application proxy (proxy firewall)}
A firewall that terminates each connection itself, rather than merely forwarding packets: it accepts
the client's connection, opens an entirely separate connection to the real destination on the
client's behalf, and relays application-layer content between the two after inspecting it.
\end{definition}

Because a proxy firewall fully reassembles and parses application-layer protocols -- HTTP requests,
FTP commands, email headers -- it can enforce rules no packet filter can even see: blocking a
disallowed HTTP method, stripping an unsafe attachment type, or refusing a malformed request before
it ever reaches the real server. The cost is performance: terminating and rebuilding every
connection is far more expensive than forwarding packets, which is why proxy firewalls are usually
deployed only in front of a small number of specific, high-value services rather than at every
network boundary.

\begin{definition}{Next-generation firewall (NGFW)}
A firewall that combines stateful inspection with deep packet inspection, application awareness
(identifying an application by behavior rather than by port number alone), and an integrated
intrusion-prevention capability, all evaluated as part of the same rule set.
\end{definition}

An NGFW answers a question a stateful firewall cannot: not just ``is port 443 open,'' but ``is the
traffic on port 443 actually the video-streaming application this policy is meant to block,''
regardless of which port it happens to use. That application awareness, plus the ability to inspect
decrypted payload content against known attack signatures, is what most modern perimeter firewalls
ship with today; ``next-generation'' describes a feature set layered on top of stateful inspection,
not a replacement for it.

\begin{misconception}
\textbf{``A stateful firewall is always more secure than a stateless one.''} Connection awareness
closes one specific class of gap. A stateful firewall configured with a rule that permits all
inbound traffic is exactly as insecure as a permissive stateless one -- the rules written still
determine the actual security posture.
\end{misconception}

\section{Firewall Rules and Access Control Lists}

A firewall's decisions are only as good as the rule list it is given.

\begin{definition}{Access control list (ACL)}
An ordered sequence of allow/deny statements, each matched against a packet's source address,
destination address, destination port, and protocol.
\end{definition}

\begin{definition}{First-match-wins}
The evaluation model in which a firewall reads an ACL top to bottom and applies the action of the
\emph{first} rule a packet matches; no later rule is ever consulted for that packet, even if a
later rule would also have matched and specified a different action.
\end{definition}

First-match-wins is the single most consequential fact about how ACLs behave: two lists containing
the identical set of rules, written in a different order, can produce entirely different security
outcomes. Every well-formed ACL ends, explicitly or implicitly, with an \textbf{implicit deny}: a
final rule dropping anything that matched nothing above it. This default-deny posture -- permit
only what is explicitly needed, reject everything else -- is the same principle underlying least
privilege in Chapter~\ref{ch:2.1}.

\begin{example}[Rule shadowing]
An administrator writes rule~1 as ``allow TCP any any port~22'' (SSH from anywhere, meant as a
temporary troubleshooting exception). Three lines later she writes rule~4 as ``deny TCP from
203.0.113.0/24 to 10.0.1.10 port~22'' (block a known-hostile range from reaching the file server
over SSH). Because rule~1 matches first, \emph{every} SSH packet -- including ones from
203.0.113.0/24 -- is allowed before the ACL ever reaches rule~4. Rule~4 is not wrong; it is
unreachable. This is \textbf{rule shadowing}: a broad rule earlier in the list makes a more specific
rule below it dead code. The fix is always about \emph{order}, never content -- the specific deny
must move above the broad allow, or the broad allow must be narrowed.
\end{example}

\begin{figure}[h]
\centering
\begin{tikzpicture}[
  dec/.style={draw=kickstartermaroon, thick, rounded corners, minimum width=2.6cm, minimum height=0.8cm, align=center, font=\scriptsize},
]
\node[dec, fill=kickstartercream] (r1) at (0,0) {Packet arrives};
\node[dec] (m1) at (3.0,0) {Matches\\rule 1?};
\node[dec] (m2) at (6.0,0) {Matches\\rule 2?};
\node[dec] (dots) at (8.6,0) {\ldots};
\node[dec, fill=kickstartermaroon!20] (deny) at (11.2,0) {Implicit\\deny-all};
\node[dec, fill=kickstartermaroon!10] (a1) at (3.0,-1.6) {Apply rule 1's\\action -- stop};
\node[dec, fill=kickstartermaroon!10] (a2) at (6.0,-1.6) {Apply rule 2's\\action -- stop};
\draw[-{Latex[length=2mm]}] (r1) -- (m1);
\draw[-{Latex[length=2mm]}] (m1) -- (m2) node[midway, above, font=\tiny] {no};
\draw[-{Latex[length=2mm]}] (m2) -- (dots) node[midway, above, font=\tiny] {no};
\draw[-{Latex[length=2mm]}] (dots) -- (deny) node[midway, above, font=\tiny] {no match};
\draw[-{Latex[length=2mm]}] (m1) -- (a1) node[midway, right, font=\tiny] {yes};
\draw[-{Latex[length=2mm]}] (m2) -- (a2) node[midway, right, font=\tiny] {yes};
\end{tikzpicture}
\caption{First-match-wins evaluation: the ACL is read top to bottom, and the first matching rule's
action is applied and evaluation stops -- no later rule is ever consulted for that packet.}
\label{fig:3.4-first-match}
\end{figure}

\begin{lstlisting}[language=Python, caption={An ACL evaluator using first-match-wins with implicit deny}]
def evaluate_acl(rules, packet):
    """rules: ordered list of dicts like {"src": str, "dst": str,
       "port": int|None, "proto": str, "action": "allow"|"deny"}.
       None for src/dst/port/proto means "any" (wildcard match).
       Returns "allow" or "deny" for the given packet."""
    def matches(rule, field, value):
        return rule[field] is None or rule[field] == value

    for rule in rules:
        if (matches(rule, "src", packet["src"]) and
                matches(rule, "dst", packet["dst"]) and
                matches(rule, "port", packet["port"]) and
                matches(rule, "proto", packet["proto"])):
            return rule["action"]
    return "deny"

def find_shadowed_rules(rules):
    """Returns indices of rules that can never fire because an earlier
    rule's conditions are a strict superset of theirs."""
    def subsumes(earlier, later):
        for field in ("src", "dst", "port", "proto"):
            if earlier[field] is not None and earlier[field] != later[field]:
                return False
        return True

    shadowed = []
    for i, rule in enumerate(rules):
        for earlier in rules[:i]:
            if subsumes(earlier, rule):
                shadowed.append(i)
                break
    return shadowed
\end{lstlisting}

The safe authoring pattern follows directly from this: place the most specific rules first --
exceptions and denies for particular hosts or ports -- then general allows, and finish with the
implicit or explicit deny-all. Auditing an existing ACL means checking each rule against every rule
listed above it. If a broader earlier rule already matches everything a later, more careful rule
would match, that later rule is shadowed. The ACL then silently does not do what its author
intended, no matter how correct that rule reads in isolation.

\begin{misconception}
\textbf{``A more specific deny rule automatically overrides a broader allow rule.''} Nothing
overrides anything in a first-match-wins engine. The only fact that matters is which rule is
physically listed first -- specificity has no special standing unless the rule order already
reflects it.
\end{misconception}

\section{Firewall Placement in the Topology}

Chapter~\ref{ch:3.3} drew the zone boundaries. This section places the firewall itself on each one,
because a rule set enforced in the wrong location protects nothing, no matter how carefully it is
written.

A \textbf{perimeter firewall} sits at the boundary between the organization's network and the
public internet, and enforces the coarsest policy in the whole design: deny everything inbound by
default, permit only the specific services the organization has chosen to expose. Behind that
perimeter, a \textbf{screened subnet} (also called a DMZ, for demilitarized zone) isolates the
public-facing servers -- a web server, a mail relay, a DNS resolver -- in their own segment,
reachable from the internet but walled off from the internal network by a second firewall boundary.
That second boundary is the entire point: if an attacker compromises the web server sitting in the
screened subnet, the screened subnet's inner firewall is what stops that compromise from reaching
payroll, engineering, or any other internal zone directly.

\textbf{Internal firewalls} extend the same logic inside the network, between zones of differing
sensitivity -- separating a finance segment from a general office segment, or a payment-processing
segment from everything else. Placing a firewall only at the perimeter assumes the perimeter never
fails; placing firewalls at internal boundaries too assumes it eventually will, and limits what a
single successful breach can reach. This is the same containment logic behind network segmentation
in Chapter~\ref{ch:3.3}, applied specifically to where an enforcement point -- not just a
boundary -- actually sits.

\section*{Review Questions}
\begin{reviewquestions}
\item An attacker sends a spoofed TCP packet with the ACK flag set to a destination port permitted
by a rule. Explain why a stateless firewall forwards this packet and why a stateful firewall,
using the same underlying rule, drops it.
\item Explain why ``stateful'' does not imply ``inspects application content,'' and give an example
of an attack that passes through a stateful firewall completely undetected as a result.
\item An ACL has, in order: (1) allow TCP any any port 22, (2) deny TCP 203.0.113.0/24 any port 22.
Explain what happens to an SSH packet from 203.0.113.5, and identify the minimal change to the ACL
that fixes the problem.
\item Explain why an ACL's implicit deny-all is described as a default-deny posture, and connect
this principle to least privilege from Chapter~\ref{ch:2.1}.
\item Two ACLs contain the exact same five rules, but administrator A orders them one way and
administrator B orders them differently. Explain how this alone can produce two different security
postures despite having identical rule content.
\item Using the first-match-wins flowchart, explain what it would mean, mechanically, for a rule
placed after the implicit deny-all to have any effect. Why does every ACL authoring guide treat
placement after deny-all as equivalent to deleting the rule?
\item An administrator wants to add a temporary, broad allow rule for troubleshooting without
risking that it silently shadows an important existing deny rule placed later in the list. Using
the safe-authoring pattern from this chapter, describe where the temporary rule should go and why.
\item A company places its only firewall at the internet perimeter, with public web and mail
servers on the same internal network segment as the finance and engineering teams. Explain what a
screened subnet (DMZ) and an internal firewall boundary would each add to this design, and why a
single perimeter firewall is not enough to contain a compromise of the public web server.
\end{reviewquestions}

\begin{answerkey}
\item A stateless firewall evaluates each packet's header in isolation with no memory of any
connection, so a spoofed ACK packet matching an allowed destination port is forwarded regardless of
whether a real session was ever established. A stateful firewall, using the identical rule, checks
the packet against its connection-tracking table first. Finding no matching entry in the
\texttt{ESTABLISHED} state for those endpoints, it drops the packet -- return traffic is only
recognized and permitted for sessions the firewall itself watched being set up through the
three-way handshake.
\answerlabel{3.4.A}
\item ``Stateful'' describes what the firewall tracks -- connection state, meaning whether a session
was legitimately established -- not what it inspects inside the payload. A stateful firewall
confirms a session is real and then permits all further traffic on it without examining the
application-layer content at all. Malware delivered inside an HTTPS stream on port 443 passes
through completely undetected as a result. The firewall correctly recognizes the connection as an
established, legitimate session, and has no visibility into what is actually being transmitted
inside it.
\answerlabel{3.4.A}
\item Because ACLs are evaluated first-match-wins, an SSH packet from 203.0.113.5 matches rule~1
(``allow TCP any any port 22'') first and is permitted immediately. The ACL never reaches rule~2, so
the more specific deny is unreachable dead code -- this is rule shadowing. The minimal fix is to
reorder the rules so the specific deny comes first: (1) deny TCP 203.0.113.0/24 any port 22, (2)
allow TCP any any port 22. The narrower, more urgent rule is then checked before the broad allow can
shadow it.
\answerlabel{3.4.B}
\item An ACL's implicit deny-all drops anything that matched no explicit rule above it, meaning
only traffic explicitly permitted is ever allowed through -- everything else is rejected by default.
This is the identical principle behind least privilege from Chapter~\ref{ch:2.1}: rather than
granting broad access and trying to block specific bad cases, both approaches start from zero
access and add back only what is specifically justified. This minimizes what an attacker can reach
if a case was simply never anticipated.
\answerlabel{3.4.B}
\item Because evaluation is first-match-wins, the physical order of otherwise identical rules
determines which rule actually governs a given packet, not the content of the rules themselves.
Suppose administrator A places a specific deny above a broad allow, and administrator B places them
in the reverse order. A's ACL blocks the traffic the deny targets, while B's ACL allows it straight
through via the broad rule before the deny is ever reached -- two firewalls enforcing the exact same
five rules produce opposite real-world outcomes purely because of sequence.
\answerlabel{3.4.B}
\item Evaluation stops the instant a packet matches any rule, and the implicit deny-all matches
literally every remaining packet by design. So no packet ever reaches a rule listed after it,
because the deny-all has already claimed the match and evaluation has already stopped. A rule
placed after deny-all is therefore mechanically identical to a rule that was never written: it can
never be reached regardless of what it says. That is exactly why it is treated as equivalent to
deletion rather than as a weaker or lower-priority rule.
\answerlabel{3.4.B}
\item The temporary allow rule should be placed \emph{after} the important deny rule it must not
shadow, not before it -- following the safe-authoring pattern of putting the most specific,
highest-stakes rules first. Placing the broad temporary allow above the deny would let it match and
terminate evaluation for every packet the deny was meant to catch, exactly as in the
rule-shadowing example. Placing it below instead ensures the deny is still checked first for any
packet it applies to, while the temporary allow still catches whatever traffic remains for the
troubleshooting window -- the flow of traffic still matches the stated policy, and tracing any given
packet against the ordered list confirms it.
\answerlabel{3.4.D}
\item A single perimeter firewall protects the boundary with the internet, but once the public web
server behind it is compromised, nothing stops the attacker from pivoting sideways to finance and
engineering -- every internal system sits on the same trusted segment. A screened subnet moves the
public web and mail servers into their own segment, behind a second firewall boundary separating
them from the internal network, so a compromise of those exposed servers does not automatically
reach anything else. An internal firewall boundary between finance and engineering adds a further
layer of the same containment: even a compromise that reaches one internal zone still has to cross
another enforced boundary to reach the other.
\answerlabel{3.4.C}
\end{answerkey}
